\subsection{IA}

\begin{itemize}
	\item Définition
	
	"Par IA, on entend, dans une acception large et ouverte, la tentative de construire des
	machines et des systèmes capables de penser et d'agir rationnellement"
	\footcite{zotero-item-584}
	
	\item Machine learning 
	"Le ML est un domaine de l'IA qui regroupe les méthodes, les processus et les techniques
	nécessaires pour rendre une machine ou un système informatique plus « intelligent »."
	\footcite{zotero-item-584}
	
	"L'idée
	à la base du ML est qu'une machine peut résoudre des problèmes de la même manière que l'homme, en apprenant à partir d’« expériences » et d'exemples, sans être explicitement programmée pour la résolution spécifique du problème. Ainsi, au lieu d'implémenter un algorithme concret de résolution de problèmes dans une machine, les données sont transmises à un algorithme générique qui, par un « processus d'apprentissage », devient
	lui-même capable de résoudre les problèmes."
	\footcite{zotero-item-584}
	
	il existe différentes méthodes de machine learning : 
	Apprentissage supervisé 
	Apprentissage non supervisé 
	Apprentissage profond et réseaux neuronaux artificiels (deep learning)
	
	Limites du machine learning : 
	"Une répartition arbitraire des systèmes est toutefois limitée dans la mesure où, d’une part, les systèmes entraînés ont souvent un domaine d'application spécifique
	et, d’autre part, leurs capacités sont fortement liées aux jeux de données utilisés. Les algorithmes entraînés avec différents jeux de données produisent des résultats
	différents."
	\footcite{zotero-item-584}
	
	"La raison pour laquelle telle ou telle fonction s'impose
	pendant le processus d'apprentissage ne peut être reconstituée que	rétrospectivement et favorise souvent des décisions curieuses de l'algorithme. L'un	des principaux problèmes réside dans le fait que les systèmes ne font pas la	distinction entre les corrélations et les causalités dans les ensembles de données.	L'exemple d'une « recherche de vaches sur des images » illustre bien cette situation	: les vaches se trouvent souvent dans des prairies ; le modèle commence alors à	ajouter « prairie » (couleur, formes, modèles d'images) au concept « vache ». L'algorithme n'a donc appris qu'une corrélation sur la base d'une probabilité statistique (vache et pré), mais n'a pas compris les relations."
	\footcite{zotero-item-584}
	
	Exemple de transcription basée sur le machine learning : Transkribus : peut aussi s'appliquer aux archives audiovisuelles. 
	
	Content-Based Image Retrieval : une des applications d la computer vision. Exemple de cas d'usage : Galt archives avec Archipanion \footcite{fewster}
	
	visual language models (VLMs) : combinaison d'algorithmes de computer vision et de LLMs pour : "create visual outputs based on input text, or textual outputs based on images or video" / "VLMs are able to do a variety of vision language tasks like	captioning and summarization, image generation, search and retrieval, and	object detection and segmentation"
	"However, there are still	challenges with these emerging models like biases, hallucinations, and overgeneralizations as well as concerns of cost and environmental impact	considering that VLMs combine two already complex AI architectures into one	even more complicated tool "
	\footcite{fewster}
	

	
	
	
\end{itemize}